\begin{table}[t]
\centering
\small
\caption{Differential agreement between vCAN and real hardware for the shared paired corpus.}
\label{tab:software-hardware-agreement}
\begin{tabular}{lrrr}
\toprule
Metric & Before & After & Total \\
\midrule
Paired case ID and patch state & 1000/1000 & 1000/1000 & 2000/2000 \\
Request payload & 1000/1000 & 1000/1000 & 2000/2000 \\
Response payload & 1000/1000 & 1000/1000 & 2000/2000 \\
Positive/negative response class & 1000/1000 & 1000/1000 & 2000/2000 \\
Negative-response code (NRC) & 1000/1000 & 1000/1000 & 2000/2000 \\
Attack-success outcome & 1000/1000 & 1000/1000 & 2000/2000 \\
Timeout outcome & 1000/1000 & 1000/1000 & 2000/2000 \\
State/readback agreement & not recorded & not recorded & requires paired readback capture \\
\bottomrule
\end{tabular}
\begin{minipage}{0.98\linewidth}\footnotesize
The comparison uses seed 20260710 and 1,000 identical case IDs/payloads per patch state (2,000 pairs total). Hardware run 1 is compared with the vCAN run. The response-class confusion matrix contains 756 positive/positive, 1,244 negative/negative, and zero off-diagonal cases. NRC agreement comprises 60 cases of 0x13, 1,184 cases of 0x31, and 756 positive responses without an NRC. Paired fuzz files do not contain read-before/read-after values, so state agreement is not inferred from the response code.
\end{minipage}
\end{table}
